\Anexo{Repositorio y datos experimentales}

El código, los datos y la memoria del trabajo están versionados en el repositorio \url{https://github.com/Jose14Andres/tfg-qwen-a100-memoria}. La Tabla~\ref{tab:anexoA:contenido} indica dónde se encuentra cada elemento que sustenta los resultados del Capítulo~\ref{chap:implementacion}.

\begin{table}[htbp]
  \centering
  \caption{Contenido del repositorio que sustenta los resultados.}
  \label{tab:anexoA:contenido}
  \footnotesize
  \begin{tabular}{|p{0.38\textwidth}|p{0.52\textwidth}|}
    \hline
    \textbf{Ruta} & \textbf{Contenido} \\
    \hline
    \texttt{experimentos/\allowbreak RESUMEN\_TESIS.md} & Cronología de las pruebas, método de medición y observaciones medidas \\
    \hline
    \texttt{experimentos/\allowbreak datos\_graficas.json} & Cada valor representado en las Figuras del Capítulo~\ref{chap:implementacion}, con el archivo del que procede; los valores derivados están marcados \\
    \hline
    \texttt{experimentos/\allowbreak generar\_graficas.py} & Script que genera las figuras a partir de los resultados registrados, sin modificarlos \\
    \hline
    \texttt{experimentos/figuras/} & Figuras en PDF y en PNG a 300~ppp \\
    \hline
    \texttt{experimentos/\allowbreak src/tfgbench/} & Cliente de medición en streaming empleado en el piloto de vLLM \\
    \hline
    \texttt{Pruebas\_Jose\_Ureta\_Tesis/} & Resultados y registros de los trabajos de Slurm \\
    \hline
  \end{tabular}
  \elaboradopor{José Andrés Ureta Granda}
  \fuente{Elaboración propia}
\end{table}

Para regenerar las figuras basta con ejecutar el script \texttt{generar\_graficas.py} en un entorno de Python que disponga de la biblioteca matplotlib, como describe \texttt{experimentos/LEEME.md}.
